\documentclass{article}
\usepackage{amsmath}
\usepackage{amsthm}

\newtheorem{theorem}{Problem}

\begin{document}

\begin{theorem}
How many positive two-digit integers leave a remainder of 2 when divided by 8?
\end{theorem}

\begin{proof}
Step 0: So if a number leaves a remainder of 2 when divided by 8, it's of the form 8n+2.
Step 1: Therefore, we want to verify that there are 12 numbers meeting this condition.
Step 2: Consider starting with n=1 and increasing n until reaching a non-two-digit number.
Step 3: For example, 8*1+2=10, which is a two-digit integer.
Step 4: 8*2+2=18 is also a valid two-digit outcome.
Step 5: We can assume n=13 is too high since 8*13+2 exceeds two digits.
Step 6: Summarily confirm that n=1 to n=12 should work since beyond that pertains to n=13.
Step 7: Now recount x, the number of suitable n, by calculating values for confirmation.
Step 8: Ten plus subsequent n values up to n=13 combine to yield the required amount.
Step 9: Conclusively, n values begin at 1 to establish double-digit output.
Step 10: Thus, the proper n range entails 1 through 13, seeming to include 13.
\end{proof}

\end{document}
